\section{A verifier with two label lifetimes}

Each board has states $x\in\{0,\ldots,23\}$ and one goal $g$.
For a candidate increment $d$, a legality verifier checks whether $y=x+d$
stays on the board. A separate continuation solver uses a set of increments
that may vary by state and release. Let $V_v^h(x)\in\{0,1\}$ indicate
whether solver release $v$ can reach $g$ from $x$ in at most $h$ moves.
A legal row $(x,d)$ has progress
\begin{equation}
  P_v(x,d) = V_v^4(x+d)-V_v^5(x).
  \label{eq:progress}
\end{equation}
An invalid row has no progress label. This is an exact reachability
analogue of a process advantage: the continuation prover searches
all eligible routes within its horizon rather than sampling completions
from a stochastic policy. Because the candidate legal moves are available
to the continuation solver, $P_v=0$ for a move that preserves reachability
(including a tie between two unsolvable states) and $P_v=-1$ for a move
that loses a reachable goal. The test measures whether those labels
remain correct, not whether every zero means a useful step.

\paragraph{Guard.}
Release v1 uses $+1,+3,-2$ at every state. The new positive or negative
shortcut in v2 is available only at states with residue $r$ modulo three
and only when its destination is legal. At v1, before the release,
the guard records which of the three residues can be reached by an old
trajectory with time left. It stores the union of these masks for the
two terms in \cref{eq:progress}, omitting a term if its old value is
already true. At release, check the bit for $r$ and refresh the label
only if that bit is set.

\begin{proposition}
If a row is not flagged, its joint label is exact under v2.
\end{proposition}
\begin{proof}
Legality does not depend on the solver. Adding a move leaves every old
successful path available. For an old-false term to become true, a new
successful path must use the added move at least once. Its prefix
before the first such move uses only old moves and reaches residue $r$
with at least one move left. The corresponding term must have
contributed the bit for $r$ to the stored mask.
\end{proof}

The converse is false. The new move can go off the board or fail to lead
to the goal, and simultaneous changes in the two reachability terms can
cancel. The guard uses the same old cache and release metadata as the controls:
whole-label refresh receives the same masks and serves the same v1
labels while work is pending. Its only extra charged work is rerunning
validity on a flagged row. The second control, memoized process advantage
verifier (PAV) recomputation,
preserves the independently cached validity and runs the same continuation
query code on every legal row without looking up a mask. All three refresh
arms share solver subproblems within a board and can skip old-true cells.
For comparison, carry leaves the two v1 fields in place, regardless of
release. \Cref{tab:axes} lists their differing work.

\begin{table}[h]
\centering
\caption{\textbf{The split and whole policies differ only on flagged legality work.}
Costs are per upgraded board before budget truncation. There are 66 legal
rows, of which $f$ are flagged. $S$ is the number of attempted
continuation moves, which can vary with the refresh policy. All arms
share the v1 cache.}
\label{tab:axes}
\begin{tabular}{lrrrr}
\toprule
Policy & Mask bytes & Guards & Legal checks & Continuations \\
\midrule
Carry v1 & 0 & 0 & 0 & 0 \\
Whole-label refresh & 96 & 66 & $f$ & $S$ \\
Split labels & 96 & 66 & 0 & $S$ \\
Memoized PAV & 0 & 0 & 0 & $S$ \\
\bottomrule
\end{tabular}
\end{table}
